%%
% ICICS / Springer LNCS（本目录 llncs.cls + splncs04.bst）。
% 正文自 paperzh short.tex 迁入并按 samplepaper.tex 版式调整（runningheads、无 geometry）。
% 编译（在 ICICS模版 目录下）：
%   pdflatex ciphernova_icics
%   bibtex ciphernova_icics
%   pdflatex ciphernova_icics
%   pdflatex ciphernova_icics
% 依赖：本目录 refs.bib（与仓库根目录 refs.bib 请自行保持同步）；插图等见 graphicspath。
\documentclass[runningheads]{llncs}

\usepackage[UTF8]{ctex}
\usepackage{graphicx}
\graphicspath{{../figures/}{../}{figures/}}
\usepackage{float}
\usepackage{amsmath,amssymb}
\usepackage[square,numbers,sort&compress]{natbib}
\usepackage{enumitem}
\setlist[enumerate]{leftmargin=1.2em, itemindent=0em, labelwidth=1em, labelsep=0.4em}
\setlist[itemize]{leftmargin=1.2em, labelsep=0.4em}
\usepackage[table]{xcolor}
\usepackage{booktabs}
\usepackage{array}
\usepackage{multirow}
\usepackage{hyperref}
\hypersetup{colorlinks=true, linkcolor=blue, citecolor=blue, urlcolor=blue}
\usepackage{caption}
\captionsetup[figure]{aboveskip=4pt, belowskip=4pt}
\captionsetup[table]{aboveskip=4pt, belowskip=4pt}
\setlength{\textfloatsep}{6pt plus 2pt minus 1pt}
\setlength{\intextsep}{6pt plus 2pt minus 1pt}

\begin{document}
\title{CIPHERNOVA：一种基于测试反馈与分层提示的密码学代码生成方法}
\titlerunning{CIPHERNOVA：测试反馈与分层提示驱动的密码学代码生成}

% ===== Anonymous switch =====
\newif\ifanonymized
\anonymizedtrue      % 投稿匿名版
% \anonymizedfalse   % 终稿实名版

\ifanonymized
\author{Anonymous Authors}
\institute{Anonymous Institution\\
\email{anonymous@anonymous.edu}}
\authorrunning{Anonymous Authors}
\else
\author{Jia-Yi He\inst{1,2} \and Xiao Lan\inst{2,1}\thanks{Corresponding author: lanxiao@scu.edu.cn}}
\institute{四川大学网络空间安全学院，成都，四川，中国\and
教育部四川大学数据防护与智能管理重点实验室，成都，四川，中国\\
\email{2022141500268@stu.scu.edu.cn, lanxiao@scu.edu.cn}}
\authorrunning{J.\ He and X.\ Lan}
\fi
\maketitle

\begin{abstract}
大语言模型（LLM）用于密码学代码生成时常面临高错误率与弱领域适配。本文提出 CIPHERNOVA 系统化方案：（1）测试反馈自改进——基于 NIST 等标准向量的闭环、错误归类与针对性修复；（2）四层可组合提示及\textbf{失败驱动提示精炼}，配套四类密码学错误分类与修复策略；（3）离线本地引擎与\textbf{教师示例检索注入}，以云端通过测试的实现为教师，在生成/改进阶段注入少样本与同题参考且无需重训练。在覆盖六种 LLM、常见分组/公钥算法及 Python/C/C++ 的 3095 条历史记录上，云端模型功能测试通过率可达 100\%；本地 Qwen2.5-Coder:7B 在更高耗时下 FTPR 亦为 100\%，完整方法在消融中显著优于无提示基线。结论限于本文基准、固定评测脚本与实验设置。
\end{abstract}

\keywords{密码学代码生成 \and 大语言模型 \and 测试反馈 \and 提示工程 \and 错误模式分类 \and 代码自改进}

\section{引言}
\subsection{研究背景与问题陈述}
密码实现须对齐分组/公钥等多形态与 NIST、FIPS 等规范；填充、初始向量等细节疏漏即可危及机密性或完整性。Copilot、CodeT5~\cite{codet5} 及主流代码大模型~\cite{codex,codellama,deepseekcoder,starcoder} 提升了通用编程辅助，但密码任务仍失败率高、依赖人工排障，且 Pass@k~\cite{codex} 等难以刻画密码代码的功能与安全诉求。为此有必要为 LLM 提供可落地的密码学代码生成优化路径。
\subsection{错误密码代码的安全影响}
\label{subsec:security_implications}
CVE 与产业案例表明，弱 IV/ECB、弱 KDF、填充错误等可致机密性或完整性失效~\cite{secsurvey,cryptollm}。本文将验收标准由\textbf{可编译}提升为\textbf{通过标准向量并满足密码学约束}，以降低 LLM 输出在安全敏感场景中的风险。
\subsection{研究目标与主要贡献}
本文提出基于测试反馈与分层提示的密码学代码生成：以闭环与分层模板引导 LLM，提高标准符合性并降低开发者负担。

主要贡献为：（1）\textbf{错误模式分类驱动的测试反馈自改进}，自动产出针对性修复指导；（2）\textbf{可组合领域提示}——四层架构与迭代精炼；（3）\textbf{离线本地引擎与教师检索注入}——无需公网即可部署；以云端通过实现为教师，在生成/改进阶段注入少样本与同题参考，\textbf{免重训练}提升本地通过率并兼顾隐私。我们在 3095 条真实用例上完成对比与消融，并提供可复现仓库。

\subsection{论文组织结构}
第 2 节综述相关工作；第 3 节给出整体架构与核心设计；第 4 节实验验证；第 5 节总结全文。

\section{相关工作}
下文从代码 LLM、领域代码生成、提示工程与自改进四方面综述。

\subsection{大语言模型代码生成}
LLM 代码生成依赖大规模预训练与自回归建模。CodeGen/ CodeGen2~\cite{codegen1,codegen2}、Codex~\cite{codex}、CodeT5/+~\cite{codet5,codet5plus}、InCoder~\cite{incoder} 及 StarCoder、WizardCoder、Code Llama、DeepSeek-Coder~\cite{starcoder,wizardcoder,codellama,deepseekcoder} 等推动了通用程序合成；CodeBERT/GraphCodeBERT~\cite{codebert,graphcodebert} 与 CodeXGLUE~\cite{codexglue} 完善了表示与评测基座。上述路线多面向通用而非密码学代码。

\subsection{领域代码生成}
领域注入有助于提升专业代码质量。EncGPT~\cite{encgpt}、CryptoLLM~\cite{cryptollm} 与 SecCoder~\cite{seccoder} 等表明领域注入有效。本文将错误类型写入提示与修复指导，形成可组合约束。

\subsection{提示工程}
CoT/Plan-and-Solve~\cite{cot,zeroshotcot,planandsolve}、少样本综合~\cite{fewshot,progsynprompt}、检索式提示~\cite{coderagbench,repolevelprompt}、分层与 Self-Refine~\cite{hierprompt,selfrefine} 改善结构与迭代，但鲜有面向密码学的专门提示。\textbf{本文}经失败驱动精炼得到\textbf{分层可组合提示}。

\subsection{代码自改进技术}
“生成—验证—改进”范式见于 Self-Edit~\cite{selfedit}、RefineCoder~\cite{refinecoder}、AgentCoder~\cite{agentcoder}、RepairAgent~\cite{repairagent}、ReCode~\cite{recode} 、学习式调试~\cite{cycle,rlef}、LLM 修复/多智能体调试~\cite{ledex,li2023,rgd,unidebugger,tang2024}、安全代码微调~\cite{secsurvey,securecodegen}。总体仍缺少面向密码学向量测试的错误分类与模板化修复。

\subsection{研究空白小结}
\label{subsec:research_gaps}
少有工作同时具备\textbf{可操作的密码学错误分类}、\textbf{随算法/语言/LLM 适配的提示}，以及与\textbf{标准向量紧耦合的自改进}。本文给出一体化方案。

\section{方法设计}
承袭通用代码 LLM 与领域注入~\cite{encgpt,seccoder}、分层/迭代提示~\cite{hierprompt,selfrefine}、闭环自改进~\cite{selfedit,refinecoder,agentcoder}，本节给出方法要点。如图~\ref{fig:framework}，\emph{CryptoAgent} 负责提示拼装、调用、提取/清洗与验证；\emph{LLMAdapter} 对接云端与本地模型；失败触发\textbf{模块 A} 闭环。系统由测试反馈自适应改进\textbf{A}、分层可组合领域提示\textbf{B}、密码学错误模式分类\textbf{C} 构成。

\begin{figure}[htbp] 
  \centering
  \includegraphics[width=\textwidth]{Framework.png} 
  \caption{基于测试反馈与分层提示的密码学代码生成系统整体框架}
  \label{fig:framework}
\end{figure}



\subsection{基于测试反馈的自适应改进机制}
传统代码生成在测试失败时仅简单重试，难以定位密码学特有的算法/实现/环境错误，通用修复难以覆盖填充、IV 等专项需求，从而导致无效加密。

\subsubsection{}该机制包含三个核心组件：错误模式识别、测试反馈分析与针对性修复。
\begin{enumerate}
    \item \textbf{错误模式识别}从输出完整性、一致性、算法完整性与占位符等维度归类，并辅以语法检查。
    \item \textbf{测试反馈分析}将测试输出压缩为缺失输出、长度差、明文误输出等结构化特征。
    \item \textbf{修复指导生成}按特征进行\textbf{模板匹配 + 规则生成}，动态填入算法、语言与错误细节。
\end{enumerate}

\subsubsection{}基于上述核心组件，本文构建了六步闭环改进机制，如图 \ref{fig:loop-process} 所示：
\begin{enumerate}
    \item \textbf{初始生成}。向 LLM 输入分层组合提示，生成初始代码；
    \item \textbf{代码验证}。以 NIST 等标准向量测试，收集输出与错误反馈；
    \item \textbf{错误分析}。通过错误模式识别生成错误类型特征向量，定位核心错误点；
    \item \textbf{修复指导生成}。根据错误特征向量输出针对性修复指令；
    \item \textbf{代码改进}。结合修复指令引导 LLM 生成改进代码；
    \item \textbf{迭代终止}。若通过测试或达到最大迭代次数则终止；否则回到第 2 步。
\end{enumerate}
\begin{figure}[htbp]
  \centering
  \includegraphics[width=0.85\textwidth]{SixStepIteration2.pdf}
  \caption{六步闭环改进流程}
  \label{fig:loop-process}
\end{figure}
% 图~\ref{fig:adaptive-iteration} 以 DES-CFB 为例，未通过时在验证—分析—改进间迭代，至多 3 轮。
% \begin{figure}[htbp]
%   \centering
%   \includegraphics[width=0.50\textwidth]{adaptive_improvement_iteration.pdf}
%   \caption{Complete iteration of the adaptive improvement mechanism for DES-CFB mode code generation}
%   \label{fig:adaptive-iteration}
% \end{figure}
\subsection{分层可组合的领域提示工程方法}
该方法旨在通过设计分层提示模板，以四层模板引导 LLM。针对代码跨场景与跨模型兼容性差等问题，该架构设计对通用规范、算法特性、模型特点与语言特性进行针对性优化。
\subsubsection{}
    \item \textbf{通用层}定义角色、规范、I/O与基线。\textbf{LLM层}按提供商与上下文做提示适配，减少跨场景错误。\textbf{语言层}按编程语言明确库调用、风格与常见错误规避，提高可编译性。\textbf{算法层}聚焦算法特有原理、步骤与参数，抬高向量验证通过率。
\subsubsection{}
    \item 该架构由最耗时的\textbf{失败迭代驱动}沉淀修复而成。实验结果的高 FTPR 依赖跨算法–语言–LLM 反复试错。\textbf{生成}侧固化头文件、自包含、按字节密文、CBC 与校验，搭配正/反例注入。\textbf{改进}侧拼接失败与编译/测试反馈，回注关键头文件、IV 与关键 API。

\subsection{面向本地模型的教师示例注入（检索式蒸馏）}
\label{subsec:distillation}
为提升离线 Qwen2.5-Coder:7B 的 FTPR，以云端通过实现为\textbf{教师}，按算法/模式/语言检索注入；改进阶段将测试反馈、修复摘要、教师块与学生代码拼入提示。首轮注入同分布少样本，后续追加同题参考；效果见表~\ref{tab:distillation_overall}。
\begin{table}[H]
  \centering
  \footnotesize
  \caption{教师示例注入前后整体通过率（Qwen2.5-Coder:7B）}
  \label{tab:distillation_overall}
  \begin{tabular}{l c c}
    \toprule
    & 通过数 & 通过率 \\
    \midrule
    注入前整体通过率 & 37 / 45 & 82.2\% \\
    注入后整体通过率 & 45 / 45 & 100\% \\
    提升             & +8 项   & +17.8\% \\
    \bottomrule
  \end{tabular}
\end{table}

\subsection{密码学代码错误模式分类系统}
\subsubsection{错误模式分类。}
本文基于真实实验归纳出四大类错误模式（如表~\ref{tab:error_classification}）。
\begin{table}[htbp]
  \centering
  \caption{密码学代码错误模式分类（四类）}
  \label{tab:error_classification}
  {\scriptsize\setlength{\tabcolsep}{6pt}
  \begin{tabular}{p{0.3\textwidth}p{0.65\textwidth}}
    \toprule
    大类 & 代表性示例 \\
    \midrule
    算法理解错误 & 模式混淆（如 OFB-64 与 OFB-8）、填充/密钥编排错误。 \\
    实现细节错误 & IV 使用（默认/随机或输出中含 IV）、填充处理、字节序。 \\
    环境配置错误 & 环境变量命名（KEY 与 TEST\_KEY）、缺库、OpenSSL 版本。 \\
    代码结构错误 & \texttt{goto cleanup} 无效、缺头文件、辅助函数不完整（未定义引用）。 \\
    \bottomrule
  \end{tabular}
  }
\end{table}

\subsubsection{修复策略设计。}
算法理解类侧重补充领域知识（模式/密钥长度/填充/库语义）；实现细节类强调 IV、字节序与输出格式；环境与结构类对齐变量名、依赖及头文件与辅助函数完整性。综合精确匹配、模式匹配与启发式规则，\textbf{错误模式匹配算法}为修复指导提供稳定标签。

\section{实验与结果分析}
% 本节通过对比、消融与错误统计实验共同验证方法有效性。
\subsection{实验设置}
\subsubsection{数据集构建。}
\label{subsec:eval_protocol}
见表~\ref{tab:dataset_stats}，在本文环境下，基于失败驱动的prompt在6种LLM上生产了3095条测试通过的真实用例，\textbf{不}等同于开放部署普适结论。
% 数据来源：仓库根目录 code_history.db，表 code_history，条件 test_success=1（与 test_data.yaml 标准向量评测一致的成功入库记录）。
\begin{table}[htbp]
  \centering
  \caption{密码学代码生成数据集统计（3095 条；按算法与语言分布）}
  \label{tab:dataset_stats}
  \setlength{\tabcolsep}{3pt}
  \begin{tabular}{p{0.12\textwidth}
    >{\centering\arraybackslash}p{0.16\textwidth}
    >{\centering\arraybackslash}p{0.30\textwidth}
    >{\centering\arraybackslash}p{0.30\textwidth}
    }
    \toprule
    算法 & 测试用例数 & 语言分布 (C/C++/Python) & 标准测试向量来源 \\
    \midrule
    DES       & 2520 & 697/693/1130 & FIPS 46-3 \\
    AES       & 335  & 102/116/117  & FIPS 197 \\
    SM4       & 184  & 60/63/61     & GM/T 0002-2012 \\
    RSA       & 56   & 19/19/18     & FIPS 186-4 \\
    \midrule
    合计      & 3095 & 878/891/1326 & -- \\
    \bottomrule
  \end{tabular}
\end{table}
\subsubsection{基线系统选择。}
为验证所提方法的有效性与相对优势，本文设置三类对照：\emph{通用单提示外部基线}、\emph{相关工作对比}、\emph{自身消融}。
\subsubsection{评估指标定义。}
为全面衡量密码学代码生成质量，本文定义六项评估指标，涵盖正确性、效率与修复能力：\textbf{生成成功率（GSR）}为成功生成可运行代码的测试用例比例；\textbf{验证通过率（VPR）}为通过语法检查与编译的测试用例比例；\textbf{功能测试通过率（FTPR）}为输出与标准测试向量一致的测试用例比例；\textbf{首次生成即通过率（FGPR）}为合并明细中 \texttt{attempts=1} 的生成会话（当次未触发会话内多轮重试）里功能测试通过的比例；\textbf{平均生成时间（AGT）}为每测试用例的平均生成时间（单位：秒）。
\subsection{对比实验结果与分析}
\subsubsection{整体性能对比。}
表~\ref{tab:llm_performance} 与图~\ref{fig:llm_performance_overview} 汇总各 LLM 上的整体表现。

\begin{table}[H]
  \centering
  \footnotesize
  \caption{所提方法在不同 LLM 上的整体性能}
  \label{tab:llm_performance}
  \setlength{\tabcolsep}{4pt}
  \begin{tabular}{
    p{0.10\textwidth}               
    >{\centering\arraybackslash}p{0.19\textwidth}  
    >{\centering\arraybackslash}p{0.20\textwidth}  
    >{\centering\arraybackslash}p{0.06\textwidth}  
    >{\centering\arraybackslash}p{0.06\textwidth}  
    >{\centering\arraybackslash}p{0.07\textwidth} 
    >{\centering\arraybackslash}p{0.07\textwidth}  
    >{\centering\arraybackslash}p{0.08\textwidth} 
  }
    \toprule
    LLM               & 提示长度 (Python/C/C++) & 生成时间 (Python/C/C++) (s) & GSR (\%) & VPR (\%) & FTPR (\%) & FGPR (\%) & AGT (s) \\
    \midrule
    OpenAI            & 17542/28952/22986            & 33.5/167.6/123.9                   & 100    & 100    & 100     & 84.80     & 92.10   \\
    Deepseek          & 8810/17227/13345             & 19.8/160.9/99.7                    & 100    & 100    & 100     & 95.73     & 42.40   \\
    Doubao            & 14075/24996/19976            & 68.8/511.8/463.5                   & 100    & 100    & 100     & 96.65     & 158.41  \\
    Claude            & 8334/15959/12139             & 28.7/153.5/166.8                   & 100    & 100    & 100     & 100       & 43.72   \\
    CodeX             & 17696/29803/23204            & 132.5/317.6/254.4                  & 100    & 100    & 100     & 100       & 117.03  \\
    qwen2.5-coder:7b  & 13972/23409/18577            & 2718.9/5647.9/4665.4               & 100    & 100    & 100     & 63.71     & 4430.14 \\
    \bottomrule
  \end{tabular}
\end{table}

\begin{figure}[H]
  \centering
  \includegraphics[width=\textwidth]{figures/llm_performance_overview.pdf}
  \caption{各 LLM 性能概览}
  \label{fig:llm_performance_overview}
\end{figure}
云端 LLM 的性能稳定居高，Claude整体性能最优；qwen2.5-coder:7b 本地大模型 FTPR 为 88.89\%，适合隐私优先场景。

\subsubsection{相关工作对比。}
见表~\ref{tab:rw_same_protocol_des}：DES 子集上，相关工作基线 FTPR 明显低于本文（DeepSeek）。

\begin{table}[H]
  \centering
  \footnotesize
  \caption{同设置相关工作对比（DES 全模式 × DeepSeek）}
  \label{tab:rw_same_protocol_des}
  \setlength{\tabcolsep}{4pt}
  \begin{tabular}{p{0.42\textwidth} c c c}
    \toprule
    基线 & GSR (\%) & VPR (\%) & FTPR (\%) \\
    \midrule
    Self-Refine+DeepSeek & 100 & 83.33 & 58.33 \\
    SecCoder+DeepSeek & 100 & 66.67 & 66.67 \\
    AgentCoder & 100 & 41.67 & 0.0 \\
    本文方法+DeepSeek & 100 & 100 & 100 \\
    \bottomrule
  \end{tabular}
\end{table}

\subsubsection{错误模式分类效果。}
对 DeepSeek 写入的功能测试失败记录经启发式归入四类；修复后合计通过率 92.80\%（表~\ref{tab:error_repair}），其中算法理解类占比最高（45.08\%），实现细节类次之。

\begin{table}[H]
  \centering
  \footnotesize
  \caption{不同错误类型的分布与修复效果（DeepSeek）}
  \label{tab:error_repair}
  \setlength{\tabcolsep}{3pt}
  \begin{tabular}{
    p{0.28\textwidth}                       
    >{\centering\arraybackslash}p{0.08\textwidth} 
    >{\centering\arraybackslash}p{0.12\textwidth} 
    >{\centering\arraybackslash}p{0.16\textwidth}  
    >{\centering\arraybackslash}p{0.18\textwidth}  
  }
    \toprule
    错误类型               & 数量 & 占比 & 修复前通过率 & 修复后通过率 \\
    \midrule
    算法理解  & 119    & 45.08\%    & 0.0\%                & 94.12\%               \\
    实现细节  & 75    & 28.41\%    & 0.0\%                & 90.67\%               \\
    环境配置  & 18    & 6.82\%    & 0.0\%                & 100\%               \\
    代码结构  & 52    & 19.70\%    & 0.0\%                & 90.38\%               \\
    \midrule
    合计      & 264   & 100\%    & 0.0\%                & 92.80\%               \\
    \bottomrule
  \end{tabular}
\end{table}

\subsection{消融实验}
消融实验用于量化所提方法各核心模块的贡献。表~\ref{tab:ablation_main} 消融表明：反馈闭环作用于通过性；分层模板不可或缺；撤除提示则指标崩溃。

\begin{table}[H]
  \centering
  \footnotesize
  \caption{主要消融实验结果（DeepSeek、豆包与 Claude）}
  \label{tab:ablation_main}
  \setlength{\tabcolsep}{4pt}
  \begin{tabular}{l c c c}
    \toprule
    系统配置                  & VPR (\%) & FTPR (\%) & 性能下降 (VPR/FTPR) \\
    \midrule
    完整所提方法                  & 97.92    & 97.92     & -- / --                     \\
    无测试反馈改进         & 100    & 81.25     & 2.08\% / -16.67\%            \\
    无分层提示架构（仅 base\_prompt.yaml） & 64.58    & 58.33     & -33.34\% / -39.59\%          \\
    无任何提示 & 39.58    & 0.0     & -58.34\% / -97.92\%         \\
    \bottomrule
  \end{tabular}
\end{table}


\begin{table}[H]
  \centering
  \caption{提示策略消融实验结果（DeepSeek、豆包与 Claude）}
  \label{tab:ablation_prompt}
  \setlength{\tabcolsep}{2pt}
  \begin{tabular}{l c c c}
    \toprule
    提示策略               & GSR (\%) & VPR (\%) & FTPR (\%) \\
    \midrule
    单一通用模板         & 100    & 45.83    & 25.0      \\
    通用+LLM              & 100    & 87.5     & 79.17     \\
    通用+LLM+语言        & 100    & 91.67    & 91.67     \\
    通用 +LLM+ 语言+算法（Proposed） & 100 & 100 & 100 \\
    \bottomrule
  \end{tabular}
\end{table}
表~\ref{tab:ablation_prompt} 显示随层数增加 VPR/FTPR 抬升，完整方法达到 FTPR 100\%。

\section{结论与未来工作}
\subsection{研究结论}
针对通用 LLM 在密码学代码上错误率高、领域适配不足的问题，本文将测试反馈闭环、四层可组合提示、\textbf{失败驱动的提示精炼}与错误模式分类整合为一体式框架。在本文实验规模下，生成/提取成功率达 100\%；完整方法性能显著优于单一通用提示。该方法在\emph{框架}上统一错误分类、闭环与分层提示；\emph{技术}上高性能提取与失败驱动精炼相配合；\emph{工程}上开源便于复现。
\subsection{研究局限}
结论依托固定评测环境与有限任务，覆盖限于常见分组/公钥算法及典型模式；评测以向量一致性为主，不宜外推为线上安全保证。
\begin{credits}
{\zihao{-5}
\subsubsection{致谢}
本研究由四川大学资助。

\subsubsection{利益声明}
作者声明在与本研究主题直接相关的范围内不存在竞争性利益。
}
\end{credits}

\renewcommand{\refname}{参考文献}
\renewcommand{\bibname}{参考文献}
\bibliographystyle{splncs04}
\bibliography{refs} 
%% Appendix (if needed)
%\appendix
%\section{Appendix Title}
%Appendix content here.

\end{document}